% !TEX root = ../main.tex
\chapter{Konklusion}
\label{ch:konklusion}

Projektet stillede spørgsmålet om, hvordan en humanoid robotarm
baseret på Energinets eksisterende 3D-design og software kan
fremstilles, samles og integreres til en funktionel enhed med korrekt
samspil mellem de mekaniske og elektriske delsystemer. Konklusionen
samler de svar, projektet har leveret, kapitel for kapitel.

\textbf{Mekanisk fremstilling og samling.} Venstre arm er printet,
efterbearbejdet og samlet til en arbejdsklar enhed. De
samlingsproblemer, der opstod undervejs -- snævre servohuse,
akselpasninger med slør, kabelgennemføringer -- er løst ved manuel
efterbearbejdning, og armen bevæger sig mekanisk frit gennem hele det
konfigurerede arbejdsrum. Detaljer i kapitel~\ref{ch:mekanik}.

\textbf{Elektrisk integration.} Armen er forbundet til en ESP32-boks
og til Jetson Orin Nano-hosten via stabil portbinding i
\enquote{/dev/serial/by-path/}. Den tidligere rækkefølgeafhængighed i
\enquote{/dev/ttyUSBx} er dermed fjernet som fejlkilde. Signal- og
effektveje er adskilte, og montagen er færdig og mærket. Detaljer i
kapitel~\ref{ch:el}.

\textbf{Software-integration og styring.} Armen kan styres
positionsmæssigt led for led via slider-GUI'en og kan afspille
præindspillede animations-sekvenser. Step-respons køres på alle
syv venstre arm-led, og resultaterne dokumenteres i bilaget med
det forbehold, der er beskrevet i afsnit~\ref{sec:hvad-maaler-vi}:
målingerne karakteriserer command-pipelinen, ikke ST3215-servoens
lukkede regulator, fordi ROS-laget bare videresender målpositioner uden at regulere selv.
Detaljer i kapitel~\ref{ch:idrift}.

\textbf{Tilføjelser til det oprindelige system.} Som del af projektet er
der udviklet en commissioning-pakke med to målerelaterede noder
(\enquote{step\_response\_node} og \enquote{repeatability\_node}) og
en samlet betjenings-GUI til opstart og demo. Disse er beskrevet i
kapitel~\ref{ch:idrift}. En per-led-effektmåling blev påbegyndt, men
er ikke leveret i fungerende stand på grund af samme
feedback-konfiguration, og den er taget ud af demoen for at undgå at
love en måling, brugeren ikke får.

\textbf{Samlet vurdering.} Problemformuleringens kerne -- at
fremstille, samle og integrere armen til en funktionel enhed med
korrekt samspil mellem mekanik og elektrik -- er besvaret. Armen
fungerer, kan styres, og de integrationsfejl, der blev fundet
undervejs, er enten løst eller eksplicit dokumenteret som vilkår,
der følger af den oprindelige arkitektur. De tilføjelser, projektet har
bidraget med ud over selve idriftsættelsen, gør det muligt at måle
og demonstrere armens opførsel og at starte den samlede system-stak
pålideligt.

